\chapter{Validación y resultados}\label{cap:validacion}

Este capítulo evalúa la corrección de la salida producida por la plataforma. Complementa el recorrido operativo \emph{end-to-end} documentado en el Anexo~\ref{ap:demo-operativa}, donde se narran las pantallas de la consola web y los artefactos producidos en cada paso. Aquí se mide qué garantiza esa salida; allí se muestra cómo se obtiene paso a paso.

\section{Estrategia de validación}

La validación combina tres niveles:

\begin{enumerate}
  \item pruebas automatizadas de módulos Python y servidor web;
  \item validación estructural del estado vivo de OpenMetadata;
  \item validación semántica del catálogo JSON-LD mediante SHACL.
\end{enumerate}

Esta estrategia evita depender de una comprobación visual de la interfaz. Cada resultado debe poder reproducirse mediante comandos versionados y artefactos generados.

\section{Pruebas automatizadas}

El repositorio contiene pruebas \texttt{pytest} para configuración, reglas de mapeo, hoja funcional, exportación DCAT, validación SHACL, validación runtime y workflow. La app web incluye tests para operaciones, lectura y escritura de CSV, resumen de jobs, artefactos y redacción de secretos. Los comandos exactos de ejecución para reproducir esta capa de pruebas figuran en el Anexo~\ref{ap:validacion-web}.

Los comandos principales se muestran en el listado~\ref{code:pytest-pnpm}.

\begin{lstlisting}[language=PowerShell, caption={Ejecución de la suite \texttt{pytest} del núcleo Python y de los tests \texttt{vitest} de la consola web.}, label=code:pytest-pnpm, captionpos=b]
$env:PYTEST_DISABLE_PLUGIN_AUTOLOAD="1"; python -m pytest
cd .\web
pnpm test
\end{lstlisting}

Estas pruebas cubren la lógica interna sin exigir que el tribunal inspeccione manualmente cada módulo.

\section{Validación estructural del estado vivo}

La validación runtime compara la instancia de OpenMetadata con el contrato esperado. Comprueba que existan servicio, base de datos, esquemas, tablas y columnas definidos por \texttt{sql/opendata\_demo\_init.sql}; también comprueba que los datasets \texttt{gold} publicados tengan los metadatos de gobierno esperados según la hoja funcional.

El comando canónico se muestra en el listado~\ref{code:validate-runtime}.

\begin{lstlisting}[language=PowerShell, caption={Validación estructural del estado vivo de OpenMetadata desde el CLI \texttt{om\_dcat\_sync}.}, label=code:validate-runtime, captionpos=b]
python -m om_dcat_sync validate-runtime --strict `
    --output tmp_pytest/runtime_validation_report.json
\end{lstlisting}

Este nivel valida trazabilidad entre la fuente técnica, OpenMetadata y la configuración funcional.

\section{Validación SHACL}

La validación SHACL usa el caso \texttt{hvd} de DCAT-AP-ES. Según la documentación oficial, este caso incluye las restricciones generales y las formas específicas para conjuntos de datos de alto valor \cite{datosGobValidacionDcatApEs2025}. En el repositorio, esas shapes están vendorizadas en \texttt{tfm\_ingestor/src/tfm\_ingestor/resources/shacl/}, junto con un manifiesto que fija el origen.

La validación puede ejecutarse de forma aislada (listado~\ref{code:validate-dcat}).

\begin{lstlisting}[language=PowerShell, caption={Validación SHACL aislada del catálogo JSON-LD ya exportado, caso HVD.}, label=code:validate-dcat, captionpos=b]
python -m om_dcat_sync validate-dcat --profile-case hvd --allow-warnings
\end{lstlisting}

La ruta recomendada, sin embargo, es integrarla dentro del workflow completo (listado~\ref{code:workflow-run}), que ejecuta sincronización, exportación y validación en un único paso idempotente.

\begin{lstlisting}[language=PowerShell, caption={Workflow canónico end-to-end: sincroniza gobierno, exporta DCAT-AP-ES y valida SHACL HVD.}, label=code:workflow-run, captionpos=b]
python -m om_dcat_sync workflow run --allow-warnings
\end{lstlisting}

\section{Validación externa con el validador europeo (SEMIC/ITB)}\label{sec:validacion-europea}

Además de la validación \ac{SHACL} local con las \emph{shapes} \ac{DCATAPES} congeladas, el catálogo exportado se ha sometido a una validación externa e independiente con el \emph{SEMIC SHACL Validator} del \emph{Interoperability Test Bed} (ITB) de la Comisión Europea~\cite{itbSemicShacl}. Este servicio comprueba la conformidad de un grafo \ac{RDF} frente a las especificaciones SEMIC oficiales. Se cargó el catálogo \fqn{validation_suite_catalog.jsonld} y se validó contra el grupo DCAT-AP, perfil \textbf{DCAT-AP HVD}, versión \textbf{3.0.0 Base Zero} (sin conocimiento de fondo), con sintaxis de entrada \ac{JSONLD}.

El resultado, recogido en la figura~\ref{fig:validacion-europea}, fue \textbf{SUCCESS} con \textbf{0 errores, 0 advertencias y 0 mensajes}. Este hallazgo es relevante por dos motivos. Primero, confirma que la salida no solo cumple el perfil nacional \ac{DCATAPES} validado localmente, sino también el perfil europeo \emph{DCAT-AP HVD} verificado por una herramienta oficial ajena a la plataforma. Segundo, la modalidad \emph{Base Zero} aplica las restricciones sin resolver vocabularios externos, de modo que una conformidad sin hallazgos evidencia que la estructura, las relaciones y las propiedades del grafo son correctas por sí mismas y no por inferencias del validador.

\figuraMemoria{fig_validation_shacl_europe.png}{Validación externa del catálogo en el \emph{SEMIC SHACL Validator} (ITB, Comisión Europea): perfil DCAT-AP HVD 3.0.0 Base Zero, resultado SUCCESS sin errores ni advertencias. Elaboración propia.}{fig:validacion-europea}

\section{Resultados del caso de uso de validación}

El caso de uso de validación produce dos datasets publicables y sus recursos asociados:

\begin{itemize}
  \item \emph{Agenda cultural pública}, a partir de \fqn{gold.agenda_cultural_publica};
  \item \emph{Resumen de movilidad por municipio}, a partir de \fqn{gold.movilidad_resumen_municipio}.
\end{itemize}

La documentación del repositorio registra una ejecución completa de la suite con los siguientes resultados: conformidad runtime, conformidad técnica, conformidad de gobierno, conformidad SHACL, idempotencia positiva, dos tablas exportadas, dos cambios aplicados en la primera ejecución y cero cambios aplicados en la segunda. Esta última cifra es relevante porque evidencia idempotencia: una vez sincronizado el estado, repetir el workflow no vuelve a aplicar operaciones innecesarias.

\begin{table}[h]
\footnotesize
\centering
\begin{tabular}{|L{5.5cm}|L{4cm}|L{4cm}|}
\hline
\rowcolor{tema!10}
\bft{Comprobación} & \bft{Resultado documentado} & \bft{Interpretación} \\ \hline
Estado runtime & Conforme & OpenMetadata contiene los activos esperados. \\ \hline
Gobierno aplicado & Conforme & La hoja funcional coincide con metadatos vivos. \\ \hline
SHACL HVD (local) & Conforme con warnings permitidos & No hay violaciones bloqueantes. \\ \hline
Validación europea (ITB) & SUCCESS (0/0/0) & Conforme con DCAT-AP HVD 3.0.0 en el validador oficial europeo. \\ \hline
Tablas exportadas & 2 & Se exportan los dos datasets \texttt{gold}. \\ \hline
Segunda ejecución & 0 cambios aplicados & El workflow es idempotente. \\ \hline
\end{tabular}
\caption{Resumen de resultados del caso de uso de validación. Elaboración propia a partir de la documentación del repositorio.}
\label{tab:resultados}
\end{table}

\section{Artefactos generados}

Los artefactos esperados son:

\begin{itemize}
  \item catálogo JSON-LD exportado;
  \item informe SHACL en Turtle;
  \item informe runtime en JSON;
  \item plan de workflow en JSON;
  \item resumen de suite de validación;
  \item informe de validación legible en HTML y PDF.
\end{itemize}

A partir del resumen JSON de la validación, la plataforma genera un \emph{informe legible} en HTML y PDF mediante el comando \texttt{om\_dcat\_sync render-report}, que se apoya en la biblioteca \texttt{fpdf2} para el PDF. El informe resalta los indicadores de conformidad (estado runtime, gobierno, idempotencia y \ac{SHACL}) y aplana el resto de resultados en un listado consultable. Puede generarse desde la consola web —pantalla \emph{Validación}— o desde el \ac{CLI}, y queda disponible como evidencia autocontenida para revisión, defensa o entrega a terceros sin necesidad de inspeccionar el JSON en bruto.

Estos ficheros se generan en \texttt{tmp\_pytest/}. No se versionan porque son reproducibles y dependen de la ejecución local, pero sí se referencian como evidencia regenerable. La cadena de pasos para regenerarlos desde un repositorio limpio se describe en el Anexo~\ref{ap:reproduccion}, y el recorrido visual que conduce a cada artefacto en el Anexo~\ref{ap:demo-operativa}.
